\section{Hardness under the Small Set Expansion Hypothesis}
\label{sec:sseh}

We prove that the Small Set Expansion Hypothesis rules out every fixed
approximation factor for treedepth. The source is the global partition
promise of Raghavendra, Steurer, and Tulsiani~\cite{RST}, in the
regular unweighted formulation of Wu, Austrin, Pitassi, and
Liu~\cite[Section~2.5]{APW}. We call
the source graph the \emph{host} to distinguish it from its line
graph. For a fixed number $q$ of classes, the YES case has a partition
into classes of equal volume, each with few crossing edges.
For fixed source parameters $c,\varepsilon>0$, the NO case gives every
set of normalized volume in $[1/10,9/10]$ at least
$c\sqrt{\varepsilon}\,m$ boundary edges. Passing to
the line graph turns the YES partition into an elimination forest of
height $O(\varepsilon m)$ and the NO cut bound into a balanced-separator
lower bound of $\Omega(\sqrt{\varepsilon}\,m)$, where $m$ is the number
of host edges and hence the order of the line graph.

Wu, Austrin, Pitassi, and Liu's incidence construction also gives
\textup{SSE} hardness for treedepth at every fixed factor when combined with
the elimination-forest bound in Appendix~\ref{app:incidence-proof}. Their output has one
part for replicated host vertices and another for host edges. Our
line-graph construction uses only host edges as vertices, joining two
when the corresponding edges share an endpoint. This gives separator
bounds directly in terms of the output order, which we use again in
the bounded-degree reductions. We prove these bounds for unweighted host
multigraphs, allowing parallel edges and loops.

\begin{theorem}\label{thm:sse-hardness}
For every fixed $C\geq1$, approximating treedepth within factor $C$ on
simple unweighted graphs is \textup{SSE}-hard.
\end{theorem}

\subsection{The source partition promise}

For a finite undirected unweighted multigraph $G$, let $m>0$ be its
number of edges. Parallel edges and loops are allowed. Edge collections
are treated as multisets, so parallel edges are counted separately.
Let $\cut_G(A)$ be the multiset of edges with one endpoint in $A$ and
the other outside $A$. The degree $\deg_G(v)$ counts the edges incident
to $v$, with a loop counted twice.
Define $\vol_G(A)=\sum_{v\in A}\deg_G(v)$ and
$\mu_G(A)=\vol_G(A)/(2m)$. When $\vol_G(A)>0$, put
$\Phi_G(A)=|\cut_G(A)|/\vol_G(A)$. Thus a loop contributes two
to the volume and never contributes to a cut.

We use the equivalent formulation of \textup{SSEH} stated in
\cite[Conjecture~2.23 and Remark~2.25]{APW}, which follows from
Raghavendra, Steurer, and Tulsiani~\cite[Theorem~3.5]{RST}.
In a regular graph, $\mu_G(A)=|A|/|V(G)|$. Taking the Gaussian error
parameter $\gamma$ to be a sufficiently small constant multiple of
$\sqrt\varepsilon$ gives the following cut bound by
\cite[Fact~2.22 and Remark~2.25]{APW}.

\begin{theorem}[Global unweighted \textup{SSE} partition gap]\label{thm:sse-source}
There are absolute constants $c>0$ and $\varepsilon_0>0$ such that,
for every fixed integer $q\geq2$ and every fixed
$0<\varepsilon<\varepsilon_0$, the following promise on regular
unweighted graphs $G$ with $m>0$ edges is \textup{SSE}-hard:
\begin{description}
\item[YES:] There is a partition
  $V(G)=V_1\mathbin{\dot\cup}\cdots\mathbin{\dot\cup}V_q$ with
  $\mu_G(V_i)=1/q$ and $\Phi_G(V_i)\leq2\varepsilon$ for every $i$.
\item[NO:] Every $A\subseteq V(G)$ with
  $\mu_G(A)\in[1/10,9/10]$ satisfies
  $|\cut_G(A)|\geq c\sqrt{\varepsilon}\,m$.
\end{description}
The source reduction has polynomial size for each fixed pair
$(q,\varepsilon)$.
\end{theorem}

The following counting lemma converts the partition properties into
bounds on crossing and internal edge counts.

\begin{lemma}[Crossing and internal edge counts]
\label{lem:sse-yes-edge-counts}
Let $G$ be a finite undirected unweighted multigraph with $m>0$ edges.
Let $(V_1,\ldots,V_q)$ be a partition of its vertices into parts
of positive volume.
Let $D$ be the multiset of edges whose endpoints lie in different
classes. For each $i$, let $E_i$ be the multiset of edges with both
endpoints in $V_i$, including loops at vertices of $V_i$. Each
parallel edge is included in the corresponding multiset. Then
$|D|\leq m\max_i\Phi_G(V_i)$ and
$|E_i|\leq\vol_G(V_i)/2$ for every $i$.
\end{lemma}

\begin{proof}
Each edge in $D$ belongs to the cuts of its two endpoint classes.
Internal edges, including loops, belong to none. Hence
\[
  2|D|=\sum_{i=1}^q |\cut_G(V_i)|
  =\sum_{i=1}^q\Phi_G(V_i)\vol_G(V_i)
  \leq\max_i\Phi_G(V_i)\sum_{i=1}^q\vol_G(V_i)
  =2m\max_i\Phi_G(V_i).
\]
Each internal edge, including a loop, contributes two to its class
volume, whereas a crossing edge contributes one. Thus
$2|E_i|+|\cut_G(V_i)|=\vol_G(V_i)$, which gives the second bound.
\end{proof}

\subsection{The line graph and its separator lower bound}

For an unweighted host multigraph $G$, its \emph{line graph}
$\mathcal{L}(G)$ has one vertex for each edge of $G$. Parallel edges
give distinct vertices. Two distinct output vertices are adjacent
if and only if their host edges share an endpoint.
This is a simple unweighted graph, even when the host has
parallel edges or loops. If the host has $m$ edges, counting parallel
edges separately, then $\mathcal{L}(G)$ has $m$ vertices.

\begin{lemma}[A treedepth upper bound from a partition]\label{lem:sse-complete}
Let $G$ be an unweighted host multigraph with $m>0$ edges, counting
parallel edges separately, and let $(V_1,\ldots,V_q)$ be a partition
of its vertices. If $D$ is the multiset of edges joining different
classes, then $\td(\mathcal{L}(G))\leq |D|+\max_i |E(G[V_i])|$.
\end{lemma}

\begin{proof}
By the vertex-integrity bound for treedepth
\cite[Section~A.3 of the full version]{GimaEtAl2022Gap}, deleting
$s$ vertices so that every remaining component has order at most
$r$ certifies treedepth at most $s+r$. Delete the line-graph
vertices corresponding to $D$. Each remaining component consists
of internal edges of a single class $V_i$, so take $s=|D|$ and
$r=\max_i|E(G[V_i])|$.
\end{proof}

For the converse direction, deleting vertices of $\mathcal L(G)$
means deleting the corresponding multiset of host edges. Different
connected components of the remaining line graph have disjoint sets of incident
host vertices. We use these sets to recover a host cut whose edges
must all have been deleted.

For $0<\rho<1$, a \emph{$\rho$-balanced separator} of an $N$-vertex
graph $H$ is a set $S\subseteq V(H)$ such that every connected component of
$H-S$ has at most $\rho N$ vertices. We write $b_\rho(H)$ for its
minimum size. A \emph{balanced separator} is a $1/2$-balanced
separator, and we abbreviate $b_{1/2}(H)$ to $b(H)$.

\begin{lemma}[From host cuts to line-graph separators]
\label{lem:linegraph-separators}
Let $G$ be an unweighted host multigraph with $m>0$ edges, counting
parallel edges separately. The vertices of $\mathcal{L}(G)$ are
indexed by these host edges, so $|V(\mathcal{L}(G))|=m$.
Suppose that, for some $\beta>0$, every set $A$ with
$\mu_G(A)\in[1/4,3/4]$ satisfies $|\cut_G(A)|\geq\beta m$.
Then $b_\rho(\mathcal{L}(G))\geq\min\{(3/4-\rho)m,\beta m\}$
for every $1/2\leq\rho<3/4$.
\end{lemma}

\begin{proof}
Fix a $\rho$-balanced separator $S$ of $\mathcal L(G)$, and write
$F_S$ for the corresponding multiset of host edges, so $|F_S|=|S|$.
If $|S|\ge(3/4-\rho)m$, the claim is immediate. Assume henceforth
that $|S|<(3/4-\rho)m$.

\emph{Grouping the surviving host edges.}
The line graph $\mathcal L(G)-S$ has
$m-|S|>(1/4+\rho)m$ vertices, with each connected component of order at
most $\rho m$. Choose a union $Z$ of connected components with
$m/4\le|Z|\le\rho m$. If a connected component has size in this interval,
take it. Otherwise, every connected component has fewer than $m/4$ vertices.
Add connected components until their total first reaches $m/4$.
The total is then below $m/2\le\rho m$.

Let $X$ be the host edges corresponding to vertices in $Z$, and
let $Y$ be the other surviving host edges. Then
$|X|=|Z|\ge m/4$ and
$|Y|=m-|S|-|X|>m/4$.
If a host vertex were incident to edges in both multisets, their
line-graph vertices would be adjacent, with one in $Z$ and one
outside $Z$. This contradicts the choice of $Z$ as a union of
connected components.

\emph{Forming the host cut.}
Let $A,B\subseteq V(G)$ be the sets of endpoints of edges in
$X,Y$, respectively. The preceding paragraph gives
$A\cap B=\emptyset$. Every edge of $X$
contributes two degree units to $\vol_G(A)$, including a loop.
The same holds for $Y$ and $B$. Thus
$\vol_G(A)\ge2|X|\ge m/2$ and
$\vol_G(B)\ge2|Y|\ge m/2$.
Since the total host volume is $2m$,
$\vol_G(A)\le2m-\vol_G(B)\le3m/2$.
Since $\mu_G(A)=\vol_G(A)/(2m)$, we obtain
$\mu_G(A)\in[1/4,3/4]$.
Every surviving edge lies wholly in $A$ or wholly in
$B\subseteq V(G)\setminus A$, so every edge crossing $A$ belongs
to $F_S$. The cut assumption therefore yields
$|S|=|F_S|\ge|\cut_G(A)|\ge\beta m$, proving the bound.
\end{proof}

\begin{lemma}\label{lem:balanced-by-td}
For every finite graph $H$, $b(H)\leq\td(H)$.
Moreover, every elimination forest $F$ of an $N$-vertex graph $H$
satisfies $\totd(F)\geq(N/2)b(H)$.
\end{lemma}

\begin{proof}
Fix an elimination forest $F$ of $H$. If every root subtree has
at most $N/2$ vertices, then $b(H)=0$ and both claims follow.
Otherwise, delete the root of the unique subtree of size greater
than $N/2$, and continue into a child subtree of size greater than
$N/2$ while one exists. Such a child is unique, since two disjoint
subtrees of that size cannot fit among $N$ vertices.

The deleted vertices form an ancestor path $P$. Every remaining
subtree has at most $N/2$ vertices, and no graph edge joins distinct
remaining subtrees. Thus $P$ is balanced and
$b(H)\leq|P|\leq h(F)$. Minimizing height proves the first claim.
The original subtree rooted at the last vertex of $P$ has more than
$N/2$ vertices, each at depth at least $|P|\geq b(H)$. Their depths
alone give $\totd(F)\geq(N/2)b(H)$.
\end{proof}

\begin{corollary}[NO treedepth bound]\label{lem:sse-sound}
Under the cut assumption of Lemma~\ref{lem:linegraph-separators},
$\td(\mathcal L(G))\geq\min\{m/4,\beta m\}$.
\end{corollary}

\begin{proof}
Take $\rho=1/2$ in Lemma~\ref{lem:linegraph-separators} and use
$b(\mathcal L(G))\leq\td(\mathcal L(G))$ from
Lemma~\ref{lem:balanced-by-td}.
\end{proof}

\subsection{The approximation gap}

\begin{proof}[Proof of Theorem~\ref{thm:sse-hardness}]
Fix $C\geq1$ and choose a fixed rational $\varepsilon>0$ small enough
that Theorem~\ref{thm:sse-source} applies,
$c\sqrt\varepsilon\leq1/4$, and $\sqrt\varepsilon<c/(8C)$.
Since $\varepsilon$ is fixed, $q=\lceil2/\varepsilon\rceil$ is a
constant. Apply that theorem with this $q$ to obtain an
unweighted host $G^*$ with $m$ edges, and output $\mathcal L(G^*)$.

In the YES case, let $D$ be the multiset of host edges joining different
partition classes. Lemma~\ref{lem:sse-yes-edge-counts} gives
$|D|\leq2\varepsilon m$ and
$|E(G^*[V_i])|\leq\vol_{G^*}(V_i)/2=m/q$.
Lemma~\ref{lem:sse-complete} therefore gives
\[
 \td(\mathcal L(G^*))\leq2\varepsilon m+m/q
 \leq(5/2)\varepsilon m<4\varepsilon m,
\]
where $q=\lceil2/\varepsilon\rceil$ gives $1/q\leq\varepsilon/2$.
In the NO case, the source cut bound holds throughout $[1/4,3/4]$.
Corollary~\ref{lem:sse-sound}, with $\beta=(c/2)\sqrt\varepsilon$, gives
\[
 \td(\mathcal L(G^*))\geq(c/2)\sqrt\varepsilon\,m
 >4C\varepsilon m,
\]
since $\sqrt\varepsilon<c/(8C)$ implies
$(c/2)\sqrt\varepsilon>4C\varepsilon$. Thus the NO bound exceeds the YES
bound by a factor greater than $C$.
\end{proof}
